\documentclass[11pt, a4paper, landscape]{article}

\usepackage[T1]{fontenc}
\usepackage[utf8]{inputenc}
\usepackage{tgpagella}
\usepackage{tgcursor}
\usepackage{microtype}
\usepackage[spanish, es-noshorthands]{babel}
\usepackage[table, xcdraw]{xcolor}
\usepackage{booktabs}
\usepackage{tabularx}
\usepackage{array}
\usepackage[most]{tcolorbox}
\usepackage[margin=1.5cm]{geometry}
\usepackage{fancyhdr}
\usepackage{amsmath, amssymb}

\pagestyle{fancy}
\fancyhf{}
\setlength{\headheight}{16pt}
\lhead{\textbf{IMT-342}}
\chead{Rúbrica Oficial — Hito 2}
\rhead{Robótica}
\lfoot{Documento de Evaluación}
\rfoot{Página \thepage}
\renewcommand{\headrulewidth}{0.4pt}
\renewcommand{\footrulewidth}{0.4pt}

\definecolor{ucbblue}{RGB}{0, 51, 102}

\begin{document}

\begin{center}
    \Large\textbf{\textcolor{ucbblue}{Rúbrica Oficial de Evaluación: Hito 2}}\\[0.2cm]
\end{center}

\section*{1. Ponderación Estructural[cite: 9]}
La calificación de la Unidad II se compone de: $\text{EC2} = 50\%$ y $\text{Hito 2} = 50\%$[cite: 9].

Dentro del Hito 2, la nota se subdivide en:
\begin{equation*}
    \boxed{ 80\% \text{ Producto Grupal} + 20\% \text{ Defensa Individual} } \text{[cite: 9]}
\end{equation*}

\section*{2. Criterios de Evaluación por Componente (80\% Grupal)[cite: 9]}

\begin{center}
\small
\renewcommand{\arraystretch}{1.3}
\begin{tabularx}{\textwidth}{|>{\hsize=0.7\hsize\raggedright\arraybackslash}X|>{\hsize=1.1\hsize\raggedright\arraybackslash}X|>{\hsize=1.1\hsize\raggedright\arraybackslash}X|>{\hsize=1.1\hsize\raggedright\arraybackslash}X|}
\hline
\rowcolor{ucbblue!10} \textbf{Dimensión (Peso)} & \textbf{Excelente} & \textbf{Parcial} & \textbf{Insuficiente} \\ \hline

\textbf{A. IK / Pieper (15\%)}[cite: 9] & 
Implementa la estructura analítica prevista; maneja correctamente ramas; usa \texttt{atan2} apropiadamente; procesa todas las poses asignadas; distingue candidatos reales/admisibles; puede explicar la geometría de las ramas[cite: 9]. & 
Obtiene resultados útiles pero existen errores menores, ramas omitidas o explicación incompleta[cite: 9]. & 
Usa un solver numérico como sustituto de la derivación; solo obtiene una solución sin justificar; no puede explicar la arquitectura de Pieper[cite: 9]. \\ \hline

\textbf{B. Validación por FK (10\%)}[cite: 9] & 
Verifica sistemáticamente candidatos; compara posición y orientación; documenta error o criterio de aceptación; identifica inconsistencias[cite: 9]. & 
Verifica solo algunos candidatos o la comparación no está bien documentada[cite: 9]. & 
Acepta soluciones de IK sin validación[cite: 9]. \\ \hline

\textbf{C. Trayectoria LSPB (15\%)}[cite: 9] & 
Construye el perfil con parámetros justificados; obtiene $q, \dot{q}, \ddot{q}$; satisface condiciones de frontera; grafica e interpreta correctamente; código y ecuaciones consistentes[cite: 9]. & 
Perfil funcional pero con validación o documentación incompleta[cite: 9]. & 
Perfil incorrecto, discontinuo o tratado como una función de biblioteca sin comprensión[cite: 9]. \\ \hline

\textbf{D. Jacobiano y Singularidades (20\%)}[cite: 9] & 
Jacobiano implementado consistentemente con al menos una validación e interpreta columnas. Usa rango como criterio central (determinante solo cuando corresponde). Desarrolla experimento de singularidad explicando la movilidad que se pierde y la amplificación de velocidades (sin confundirla con torque infinito)[cite: 9]. & 
Detección numérica correcta con interpretación física incompleta[cite: 9]. & 
Declara singularidad únicamente por observación; utiliza mal el determinante; no relaciona el resultado con la pérdida de movilidad local[cite: 9]. \\ \hline

\textbf{E. Dinámica y Torque (15\%)}[cite: 9] & 
Utiliza el modelo 1R trabajado; documenta supuestos y parámetros; calcula torque a partir de $q, \dot{q}, \ddot{q}$; separa e interpreta contribuciones; presenta $\tau(t)$; relaciona el comportamiento con la trayectoria[cite: 9]. & 
Cálculos correctos pero con poca interpretación o validación[cite: 9]. & 
Torque obtenido sin modelo claro; parámetros sin procedencia; confusión entre singularidad y torque[cite: 9]. \\ \hline

\textbf{F. Integración y Reproducibilidad (5\%)}[cite: 9] & 
El \texttt{.ipynb} ejecuta de inicio a fin; módulos comparten datos coherentemente; código legible; unidades/convenciones declaradas; informe IEEE sintetiza correctamente; resultados reproducibles[cite: 9]. & 
Integración funcional con algunas dependencias manuales o documentación insuficiente[cite: 9]. & 
Scripts aislados; notebook no reproducible; resultados no pueden reconstruirse[cite: 9]. \\ \hline

\end{tabularx}
\end{center}

\newpage

\section*{3. Criterios de Defensa Individual (20\%)[cite: 9]}
La evaluación individual durante la defensa se desglosa en cuatro subdimensiones equitativas:

\begin{center}
\small
\renewcommand{\arraystretch}{1.5}
\begin{tabularx}{\textwidth}{|>{\hsize=0.45\hsize\raggedright\arraybackslash}X|c|>{\hsize=1.55\hsize\raggedright\arraybackslash}X|}
\hline
\rowcolor{gray!20} \textbf{Subdimensión} & \textbf{Peso} & \textbf{Requisito Demostrable} \\ \hline
\textbf{G1. Pipeline}[cite: 9] & $5\%$ & Explica claramente cómo se conectan los componentes lógicos del proyecto[cite: 9]. \\ \hline
\textbf{G2. Matemática}[cite: 9] & $5\%$ & Reconstruye o explica al menos una ecuación relevante del modelo[cite: 9]. \\ \hline
\textbf{G3. Validación / Interpretación}[cite: 9] & $5\%$ & Explica cómo sabe que un resultado numérico es correcto y qué significa físicamente en el contexto del robot[cite: 9]. \\ \hline
\textbf{G4. Transferencia}[cite: 9] & $5\%$ & Responde exitosamente sobre una parte distinta del pipeline a la de su contribución principal de programación[cite: 9]. \\ \hline
\end{tabularx}
\end{center}

\vspace{0.5cm}
\noindent
\begin{minipage}[t]{0.48\textwidth}
    \begin{tcolorbox}[colback=red!5, colframe=red!70!black, title=\textbf{Regla Individual Mínima}, left=1mm, right=1mm, top=1mm, bottom=1mm]
        \small
        La evaluación individual es \textbf{obligatoria}[cite: 9]. 

        Un estudiante que no pueda demostrar una comprensión mínima del trabajo presentado podrá recibir una reducción individual directa en la nota de su Hito 2, \textbf{aun cuando el producto grupal sea técnicamente correcto}[cite: 9]. 

        La defensa está diseñada para verificar comprensión transversal, no memoria literal[cite: 9].
    \end{tcolorbox}
\end{minipage}%
\hfill
\begin{minipage}[t]{0.48\textwidth}
    \begin{tcolorbox}[colback=yellow!10, colframe=orange, title=\textbf{Políticas de Uso de IA}, left=1mm, right=1mm, top=1mm, bottom=1mm]
        \small
        El uso de Inteligencia Artificial está \textbf{permitido}[cite: 9]. Sin embargo, la evaluación recaerá estrictamente sobre la capacidad del estudiante para:
        \begin{itemize}
            \item Verificar y razonar los resultados[cite: 9].
            \item Garantizar la reproducibilidad[cite: 9].
            \item Explicar y diagnosticar el trabajo[cite: 9].
        \end{itemize}
        La respuesta pasiva \textit{“lo generó la IA”} no constituye en ningún caso una justificación técnica válida ante el tribunal[cite: 9].
    \end{tcolorbox}
\end{minipage}

\end{document}
